\cotefile{157-4}
\scanpage{379}
\mspage{\unsure{153|103}}
% archive sheet number: 378
flèches de la forme $f\times\mathrm{id}_X:S\times X\longrightarrow T\times X$, où $f:S\unsure{\relbar|\longrightarrow}T$ est un morphisme entre ensembles ordonnés finis. Ou plus généralement, entre ensembles ordonnés $S,T$ tels que $S^{\unsure{o|0}}$ soit artinien, \unsure{i.e.} $S$ noethérien, et que $\mathrm{Esp}(S^o)$ soit noethérien. Cette dernière hypothèse étant inutile sous l'axiome $\mathrm{Loc}(8)$ \unsure{des limites|en limite}, cf $\underline{\overline{\mathrm{XVII}}}$ th.~3.1 p.~56.) Donc ce n'est pas encore le résultat rêvé !

7.5. Le seul cas que je connaisse où $f_*$ existe pour \uline{toute} flèche dans $\mathrm{Cat}$, c'est le cas où $W$ satisfait :
\[
\underline{\mathrm{Loc}(\deleted{\unsure{1\ \mathrm{bis},|1,}}\ 6a,8,9)}\,,
\]
% slanted
\margin{vérifier s'il ne faut pas aussi la condition $\mathrm{Loc}(3\,\mathrm{bis})$ en plus de $\mathrm{Loc}(6a,8,9)$ \unsure{($\alpha$)} Vérification faite, c'est OK \unsure{sans} $\mathrm{Loc}(3\,\mathrm{bis})$}
\margin{\uline{NB} Après nouvelle réflexion, il \unsure{apparait|apparaît} qu'il suffit même $\mathrm{Loc}(6'a,8,9)$}
c'est à dire \deleted{l'axiome $W=W^{\unsure{(6)|6}}$ des facteurs directs,} l'axiome des $W$-asphérisations faibles, déjà abondamment mentionné, et enfin \unsure{ceux|aux} des limites et d'accessibilité, i.e. $W$ est \unsure{$L_1$|Lim}-accessible dans $\underline{\mathrm{Fl}}(\mathrm{Cat})$.

Le résultat-clef pertinent est donné dans XV 8.4. (p.~105-\unsure{107|106}), qu'il faut souder à XV 7.4. (p.~98-99) pour en faire
